\uuid{U2jp}
\titre{Congruences}
\auteur{bodin ; croizat ; sacre}
\organisation{exo7}
\licence{CC BY-NC-SA 4.0}
\datecreate{2022-03-08}
\niveau{L1}
\module{Algèbre}
\chapitre{Arithmétique dans Z}
\sousChapitre{Autre}
\difficulte{1}
\theme{Congruences}
\exo7qcmid{hilisit.arithmetique.042}

\contenu{
\begin{question}
\qclassification{103.99}
 Si $x \equiv 2 \; [5]$, alors on a :
    \begin{answers}
        \good{$x^2 \equiv 2x \; [5]$}
        \bad{$3x \equiv -1 \; [5]$}
        \good{$x+1 \equiv 3 \; [5]$}
        \bad{$10x \equiv 2 \; [5]$}
    \end{answers}
    \begin{explanations}
    D'après les propriétés arithmétiques des congruences et notre congruence initiale $x \equiv 2 \;[5]$ :\\
    en ajoutant $1$ : $x+1 \equiv 2+1 = 3 \; [5]$,\\
    en multipliant par $3$ : $3x \equiv 3 \times 2 = 6 \equiv 1 \;[5]$,\\
    en multipliant par $10$ : $10x \equiv 10 \times 2 = 20 \equiv 0 \;[5]$,\\
    Enfin on calcule : $x^2 \equiv 2 \times 2 \equiv 2x \;[5]$.
    \end{explanations}
\end{question}
}
